% ============================================================
%  part4/ch23-stability.tex  —  Chapter 23: Stability
% ============================================================

\chapter{Stability Analysis}
\label{ch:stability}

\section{Linear stability of the lamphron subsystem}

Let $L = L_0 + \delta L$ near a uniform equilibrium.
Assume $\mu_L = aL - b\Delta L$ (linearised chemical potential).

\begin{proposition}[Lamphron linear stability]
  \label{prop:lamphron-stability}
  \statusproven{}
  Fourier mode $\delta L = \hat{L}_k e^{ik\cdot x + \sigma(k)t}$
  has growth rate $\sigma(k) = -M_L(ak^2 + bk^4)$.
\end{proposition}

\begin{proof}
  Substitute into $\partial_t(\delta L) = M_L\Delta(a\delta L - b\Delta\delta L)$.
  In Fourier space: $\sigma(k) = -M_L(ak^2 + bk^4)$.
\end{proof}

\begin{proposition}[High-frequency stabilisation]
  \label{prop:hf-stable}
  \statusproven{}
  For $b > 0$ and large $k$: $\sigma(k) \sim -M_L bk^4 < 0$.
  High-frequency perturbations decay.
\end{proposition}

\begin{proof}
  For $k \gg \sqrt{a/b}$, the $bk^4$ term dominates.
  Since $M_L, b > 0$: $\sigma(k) < 0$.
\end{proof}

This establishes finite characteristic scales for lamphron
domains — a prerequisite for the cosmic structure discussion
in Chapter~\ref{ch:filament-formation}.

\begin{HolonomyNote}
  After traversing Part~IV, what is visible in retrospect?

  Chapters~19--23 look like PDE theory.
  But the structure they establish is exactly the framework
  that Part~V needs: Sobolev categories license the sheaf,
  the energy dissipation theorems license the gradient-flow
  interpretation, and the stability results license the
  claim that RSVP structures are long-lived rather than
  transient.

  The holonomy of Part~IV is: analytic regularity is not a
  technical detail.
  It is the condition of possibility for geometric structure.
\end{HolonomyNote}
